





%´´´´´´´´´´´´´´´NOW COVERED IN 07´´´´´´´´´´´´´´´´´´´´´´´´´´

%!!!!!!!!!!!!!!!!!!!!!!!!!!!!!!!!!!!!!!!!!!!!!!!!!!!!!!!!




































%!TEX root=../main.tex
\chapter{Limitation \& Future Work}%todo change title to match contents
\label{chap:limitation}
\acresetall %reset abbrevations

\begin{comment}

%tim aus meeting:
%ausblick wohin? Kapitel Limitations und Future Work
%nur die sationen, nur 1 jahe.. alle annahmen. 
%transformer..
%domänenanalyse machen, wie unterschiedlich die sind
%feature weights (abaltio bei auswertung oder mischen) pro modell -> wo ist was am wichtigsten? dadurch fazit domäne relevant
%domnenentscheidung datanbasiert, evtl statdesse nclustern etc
%vq transformer
%augmentation




\end{comment}
\todo{limitation schiff muss wissen wo es ist -> OSM interface muss bestehen; außerdem mögliche analyse wie viel data hilft wirklich, wie schnell kovergiert es wenn man sampled kann massiv zeit sparen; außerdem, consider waterboundaries if we use osm dataa anyways}

%\todo{manual sampling before running the TPE-> biased to that region, lack of exploitation-> either also add more random runs, or do HP selection befroehand, or just add many random runs that will cover all hp sets. at least the initial runs were very likely globally good, as i optimised them on a smaller set, so this can also be seen as a prior we want to bias to, though that was not my intention}

\ws{This chapter can be a section* in the conclusions chapter. No need to give it a separate chapter.}
\todo{ 
There are also notes regarding the combination of Chapters 7, 8, 9 in one. 
    I can only read these after the restructuring. Now, it won't be as effective to read the text out of correct context.
}\\
\section{Dataset limitations}

\paragraph{Dataset Coverage and AIS Biases.}
The results in this thesis are based on AIS logs from three German measurement stations and a single calendar year. While the datasets cover a large number of trajectories, they still reflect the traffic patterns of a limited set of waterways and regulatory conditions. For the observation period of the year 2025, the data availability varies across the three measurement sites: Kiel (357 days), Wedel (351 days), and Bremerhaven (326 days). In AIS networks,  data loss can occur for diverse reasons, such as hardware maintenance. Notably, there were no data for June 2025 in Bremerhaven. The gaps are not a concern in practice. Since the deep learning models operate on short, localised trajectory windows rather than continuous long-term sequences, full calendar coverage is not required. Each available day yields a large number of independent training samples, and together they cover a representative variety of traffic pattern. Extending the analysis to additional years and measurement sites systems would be necessary to assess whether the observed error distributions and monthly trends generalise well. 

Furthermore, the data primarily reflects the movement patterns of large commercial vessels, which are mandated to carry Class A transponders under the \ac{SOLAS} \parencite{solas2014}. This includes cargo vessels exceeding 300GT on international voyages and all tankers. Consequently, smaller recreational craft and fishing vessels with Class B transponders are significantly underrepresented. Because Class B devices transmit at lower power and with less frequent update rates (as detailed in Table \ref{tab:ais-frequencies}), the resulting models may be biased to the more predictable dynamics of large commercial ships.

\paragraph{Augmentation and class imbalance.}
The geometric augmentation detailed in Section \ref{subsubsec:augmentation} addresses the class imbalance for lock manoeuvres, but it also introduces specific assumptions into the training data. 
The spatial coordinates ($x, y$) and angles are transformed, but the temporal spacing ($\Delta t$) and the velocity magnitude (\ac{SOG}) remain unchanged. Thus, the model observes the exact same acceleration and deceleration patterns multiple times, only from different directions. This repetition increases the risk of the model memorizing specific kinematic profiles rather than learning generalised lock approach behaviour and potentially elevates the epistemic uncertainty when the model encounters new velocity profiles during inference.\\
Additionally, random rotation assumes spatial invariance, implicitly treating the vessel motion as independent of environmental directionality. For locks, this assumption can safely be made, but one should be careful when augmenting e.g. rivers, where local currents, or wind conditions influence manoeuvring dynamics.\\
Similarly, the mirroring operation assumes perfect lateral kinematic symmetry, neglecting real-world asymmetries such as the propeller walk , which is common in single-screw inland vessels at low speeds. While these simplifications are necessary to generate a balanced dataset, they limit the exposure of the model to direction-specific environmental disturbances.

Because only data within the lock domain were augmented, hyperparameters are distinct and hard to compare to the other domains. Future work should explore alternative strategies for handling rare domains, such as downsampling of other classes, and quantify how strongly these design choices affect the error distribution and the robustness of the models


\paragraph{Derived kinematic features.}
\label{subsubsec:kinematiclowvel}
A notable limitation in the data preprocessing pipeline stems from the fact that the \ac{ROT} is not directly available in most raw AIS messages. It was derived analytically as the time derivative of the \ac{COG}. Since \ac{COG} is calculated from consecutive \ac{GNSS} position fixes rather than being measured by a dedicated directional sensor, it can become highly unreliable at near-stationary speeds. At low velocities, the localisation noise of the \ac{GNSS} receiver exceeds the actual distance travelled. This causes the computed \ac{COG} vector to fluctuate, resulting in artificial, physically unrealistic ROT spikes that do not reflect actual vessel movement. 

While the true heading channel of the vessel offers a more stable orientation metric at low speeds, we do not use it as as an alternative feature for two primary reasons. First, as per \textcite{itu-r-m1371-1-2001}, true heading is an optional AIS field and is frequently missing, particularly among smaller inland vessels equipped with Class B transponders. Second, heading only indicates the direction the bow of the vessel is pointing, whereas \ac{COG} represents the actual path of motion over the earth. Because inland waterway vessels can be subjected to lateral drift from river currents and wind, \ac{COG} was used for predictions, despite the % \ac{GNSS}-induced 
limitations at near-zero velocities. %unterschied COG und true heading wäre interessant für den schwimmwinkel als indikator für strömunbg. geht hier aber zu weit und nur möglich wenn wir stark filtern da nicht oft verfügbar
Since both the EKF and the LSTM optimisation indicate that the derived rate-of-turn signal does not contribute much predictive value, this choice should be revisited in future works.










\section{Methodological limitations}


\paragraph{Hyperparameter optimisation bias.}
We initialised the Optuna search for the \ac{LSTM} with a manually selected subset of starting configurations (16 out of 26 seeding trials). This ensures that all feature sets are evaluated at least once, but it also introduces an implicit prior over the search space, because the \ac{TPE} sampler conditions its density estimates on these seedings. Future work should therefore compare this seeded strategy against a larger number of fully random initial trials or split up the search into two stages, to separate feature selection from the remaining hyperparameters, and quantify the difference to our implementation.


\paragraph{Statistical analysis.}
The statistical analysis in Chapter \ref{par:statistical} is based on file-level aggregated metrics, with one file representing one domain on one day on one measurement station. This is used as a fallback to per-trajectory evaluation, since only the aggregated errors were available at evaluation time. As a result, all files contribute equally as blocks, even though they may contain different numbers of trajectories. %If the Friedman test is significant , we apply the Nemenyi post-hoc test for pairwise comparisons and visualise the results using a critical difference diagram. 
Future iterations of this work should save the per-trajectory evaluation for more accurate metrics.


\paragraph{Domain labels and context representation.}
The work relies on a four-class domain labelling scheme (harbour, river, channel, lock) obtained from OpenStreetMap. %For trajectories with multiple domains, a majority vote is applied.
A natural next step would be to refine the domain representation, for example by introducing additional domain types, or by clustering waterway segments based on kinematic statistics to obtain a more fine-grained context. It would also be interesting to investigate whether representing the domain as a sequence over time, instead of a single static vector per sample, further improves performance for longer context windows.

\paragraph{Hyperparameter optimisation budget.}
Because of the pruning and early stopping mechanisms we used, hyperparameter search can terminate before the true optimum is explored. The observation that the domain-aware LSTM with reused hyperparameters outperforms the separately tuned variant suggests that at least one Optuna study stopped prematurely. Future iterations of this work should consider larger budgets or more conservative pruning schedules for the most promising architectures.



\section{Model-specific limitations}

\paragraph{LSTM architecture.}
The LSTM model implemented in this thesis uses a stacked encoder with a linear decoder that predicts the full horizon in a single step. This choice matches the short 5-minute prediction horizon, but it also restricts the flexibility of the architecture. Autoregressive decoders, where each predicted step is fed back into the model as additional context, are the standard for \acp{LSTM}. In the current setup, both context and prediction length are fixed and identical across all models, so it remains unclear how sensitive the observed ranking between models is to these design choices. Future work should therefore investigate autoregressive decoders and variable context and prediction windows.
Additionally, a bidirectional encoder should be investigated, as it has improved model performance in other scenarios \parencite{articlebilstmsus}. Other gains could be found by introducing attention mechanisms, that give higher weights to relevant timestamps of the context window. 
In addition to time-wise attention, covariate-informed forecasting with attention across feature channels is particularly relevant. The comparison between Chronos-2 and TiRex suggests that cross-channel attention can provide an advantage over univariate or channel-independent sequence models. At the moment, TiRex is limited to univariate inputs, and the LSTM uses only a simple concatenation of features without explicit channel weighting. Ongoing work on extending x-LSTM-based models such as TiRex to handle multivariate, covariate-informed inputs would offer a direct opportunity to revisit this comparison once such models become available. Beyond recurrent architectures, Transformer models for time series have shown promising results in other forecasting domains, mainly due to their ability to model long-range dependencies. Training such models from scratch on the AIS data was out of scope for this thesis because of the associated runtime and tuning cost, but remains an interesting direction for future work. \me{To Tim: Ich finde den sprung zu VQ transformer immer noch nicht sinnvoll und empfehle ihn hier daher nicht. Transformer haben eine längere Trainingszeit als LSTM. Der VQ Transformer wirkt dem entgegen für lange kontexte. Diese liegen hier aber nicht vor.}

\paragraph{Uncertainty quantification for learned models.}
 For applications in multi-target tracking and autonomous navigation, a well-calibrated estimate of predictive uncertainty is important. A natural extension would therefore be to equip the LSTM with probabilistic outputs%, for example by ... %predicting multiple quantiles, training small ensembles, or applying conformal prediction on top of the existing point forecasts.
 %\todo{look into it, types of uncertainty} Such methods 
 . That would make it possible to evaluate the calibration of LSTM-based motion models in the same way as for the probabilistic baselines and to study how uncertainty behaviour differs between domains.

%-------------------------------postambel
%\me{das hier runter sehr unwihtig}
%   - What worked, what didn't, why
%   - Generalisation, runtime, practical constraints
 %  - Integration considerations for IRT pipeline


  


%\section{Methodological Challenges and Data Anomalies}
%\me{eher für rückfragen as so challenges waren, aber da platz sowieso knapp ist ist das egal}
%Developing and migrating the data pipeline presented a few technical challenges that required specific workarounds. Early in the prototyping phase, while using a small, locally subsampled dataset, the model consistently achieved a lower loss on the validation set than on the training set. This unexpected behaviour was caused by an uneven class distribution within the small sample size. A standard random split of a highly diverse dataset created a validation set filled with simple, easy-to-predict trajectories. To fix this, a data stratification strategy was implemented.Once this stratification was applied to the full dataset the data distribution stabilised and the validation metrics behaved as expected. %wurde ja auch garnicht behoben

%Another critical issue occurred when moving the pipeline between different working environments. Running the code on a new machine led to a massive, drop in valid harbor trajectories. This was traced back to a version mismatch in the \textit{pandas} library. Between major versions, the default parsing resolution for \texttt{datetime} objects changed from nanoseconds to seconds. Because the spatial jump-detection filter relied on implicitly converting these timestamps to integers, the newer library version caused the code to calculate microscopic time-deltas ($\Delta t$). This math error made the filter compute physically impossible speeds, causing it to aggressively delete thousands of valid, slow-moving harbor maneuvers. The issue was resolved by rewriting the filter to explicitly calculate the time difference in seconds in order to work with all pandas version.

%A third issue emerged during the first full-dataset evaluation with the baseline model. The RMSE was several orders of magnitude larger than the ADE, indicating the presence of extreme coordinate outliers. These were traced to overshoot artefacts introduced by cubic spline interpolation during resampling: rarelz, the spline solver produced physically impossible position jumps of up to tens of kilometres after the jump filter had been applied. The interpolation method was replaced with PCHIP, which guarantees monotonicity within each interval and cannot overshoot the local value range. Additionally, the time-gap guard in the track segmentation step was tightened to reject any window whose actual elapsed time exceeded the theoretical context-plus-prediction duration, eliminating tracks with hidden data gaps that would otherwise produce misleading ground-truth positions.